\begin{abstract}
Ghép cặp hồ sơ ứng viên (Curriculum Vitae, CV) với mô tả công việc (Job Description, JD) là một bài toán cốt lõi trong các hệ thống tuyển dụng số và nền tảng kinh tế việc làm tự do (freelance). Các hệ thống truyền thống chủ yếu dựa trên đối sánh từ khóa từ vựng hoặc nén toàn bộ tài liệu thành một vector đặc đơn lẻ (Single-Vector Dense). Tuy nhiên, cách tiếp cận đơn vector làm loãng ngữ nghĩa của các kỹ năng kỹ thuật cốt lõi trong các đoạn văn dài, đồng thời gặp giới hạn chiều dài ngữ cảnh ($256$--$512$ tokens) của mô hình Transformer, dẫn đến cắt xén thông tin nghiêm trọng trên phần lớn CV. Ngoài ra, tuyển dụng thực tế đòi hỏi các ràng buộc nghiệp vụ cứng (kỹ năng bắt buộc so với kỹ năng ưu tiên) mà khoảng cách ngữ nghĩa thuần túy không thể đáp ứng. Trong bài báo này, chúng tôi đề xuất một kiến trúc ghép cặp hai chiều có khả năng triển khai thực tế trong công nghiệp, dựa trên biểu diễn đa vector nhận biết phân mục (Section-Aware Multi-Vector) kết hợp động cơ xếp hạng lai (Hybrid Ranking Engine). Phương pháp bóc tách tài liệu thành các phân mục chức năng độc lập (kỹ năng, kinh nghiệm, học vấn, tóm tắt) thông qua bộ phân tích quy tắc tất định không rò rỉ, mã hóa độc lập qua mô hình Bi-Encoder đóng băng, và ghép cặp tương thích chéo với cơ chế tái chuẩn hóa trọng số động nhằm xử lý công bằng các hồ sơ thiếu mục tự chọn. Điểm số tương đồng ngữ nghĩa sau đó được kết hợp cùng chỉ số từ vựng BM25 và độ phủ kỹ năng bất đối xứng trên cây phân loại (Taxonomy) 3 tầng gồm 615 kỹ năng chuẩn hóa. Thử nghiệm thực nghiệm diện rộng trên tập dữ liệu chuẩn quốc tế \textsc{resume-job-description-fit} ($1.759$ cặp nhãn trên $71$ truy vấn JD) cho thấy mô hình lai đạt MAP $0.6366$ trên toàn bộ truy vấn ($0.7410$ trên các truy vấn có ứng viên phù hợp), nDCG@10 đạt $0.6162$, và hệ số tương quan hạng Spearman $\rho_{\text{mac}} = 0.2954$. Về hiệu năng tính toán, độ trễ xếp hạng trực tuyến cho tập $N=25$ ứng viên chỉ tiêu tốn $T_{\text{rank}} = 2.69$\,ms trên CPU phổ thông ($0.12$\,ms đối với Section-Aware Dense), với tổng độ trễ trực tuyến toàn trình $T_{\text{online}} \approx 33.4$\,ms bao gồm cả mã hóa truy vấn---nhanh hơn gấp $35$ lần so với mô hình chú ý chéo Cross-Encoder ($1.178$\,ms). Kiểm định thống kê ghép cặp với $10.000$ mẫu lặp Bootstrap 95\% CI và hiệu chỉnh đa giả thuyết Holm--Bonferroni khẳng định sự tương đương thống kê về độ chính xác giữa Đa Vector và Đơn Vector phẳng trên tập test và tập huấn luyện độc lập (76 truy vấn sạch, $1.829$ cặp), qua đó làm sáng tỏ vai trò thực tế của kiến trúc nhận biết mục: triệt tiêu hiện tượng cắt xén văn bản ($>95\%$ tránh truncation) và cung cấp khả năng giải thích minh bạch hỗ trợ nhà tuyển dụng.

\keywords{Truy xuất thông tin \and Biểu diễn đa vector \and Nhận biết phân mục \and Taxonomy kỹ năng \and Xếp hạng lai \and Ghép cặp CV--JD \and Hệ thống gợi ý.}
\end{abstract}
